\documentclass[11pt,a4paper]{article}
\usepackage[margin=0.85in,top=1in,bottom=1in]{geometry}
\usepackage{amsmath,amssymb,amsfonts}
\usepackage{graphicx}
\usepackage{float}
\usepackage{enumitem}
\usepackage{microtype}
\usepackage{caption}
\usepackage[hidelinks]{hyperref}
\usepackage[most,breakable]{tcolorbox}
\tcbuselibrary{breakable,skins}
\usepackage[utf8]{inputenc}
\usepackage[T1]{fontenc}
\usepackage{lmodern}
\captionsetup{font=small,labelfont=bf,skip=4pt}
\setlist{leftmargin=1.5em,topsep=2pt}

%% Breakable card — prevents content cutoff across pages
\newtcolorbox{solcard}[3]{
  breakable, enhanced,
  colback=white, colframe=gray!50, arc=2pt,
  title={\textbf{#1}\hfill\textsf{\small #3\,pts}},
  fonttitle=\normalsize\sffamily\bfseries,
  before upper={\small\textit{#2}\par\smallskip\hrule height 0.4pt\smallskip},
  top=4pt, bottom=6pt, left=7pt, right=7pt,
  parbox=false,
}

\title{\textbf{Rolling and sliding friction --- Solution}}
\author{\normalsize W21 (2021) — English translation}
\date{}
\begin{document}\maketitle\vspace{-0.5em}
\begin{solcard}{A1}{Determine experimentally the rolling friction coefficient $k$ of a cylinder rolling on a wooden block. Describe the measurement method and sketch the installation diagram.}{11.00}

Let's balance the ruler on the edge of the block, thereby determining the position of the center of mass of the ruler: $x_r=248~\mathrm{mm}$. Using the ruler as a lever, we determine the mass ratio $M/m$ of the cylinder and the ruler: $$M/m=3.44 \pm 0.22.$$ 

Attach the ruler as a rod to the base of the cylinder using an adhesive mass. Direct the side of the ruler with the scale towards the cylinder: this makes it convenient to determine the ruler's position relative to the cylinder. Attach a sheet of paper to the edge of the table and draw lines at different angles to the vertical. Set the cylinder and rod system into oscillation. Record the amplitude values by aligning them with the lines on the sheet. The lines on the sheet do not necessarily need to be drawn as rays from a single center: one can select positions so that the lines coincide with the inclined ruler for the corresponding positions of the cylinder. 

The energy of the oscillating system, expressed in terms of the angular amplitude $\varphi$: $$W=mgd(1-\cos\varphi ) \approx mgd \frac{\varphi^2}{2}.$$ 
For a rotation of the cylinder with the rod by an angle $d\alpha$, the torque of the reaction force performs work: $$\delta A=-M \cdot d\alpha=-k(m+M)g \cdot d\alpha.$$ 
The remaining forces do not perform work. Note that even when the system passes through the equilibrium position and the motion has not yet ended, the system is still characterized by the angular amplitude $\varphi$. The change in the system's energy for a rotation by a small angle $d\alpha$ is: $$dW=mgd \cdot \varphi d\varphi=-k(m+M)g \cdot d\alpha.$$ 
The decrease in energy over 1/4 of a period (here the symbol $\Delta$ means decrease, not change) is: $$mgd \cdot \varphi \cdot \Delta\varphi_{T/4} = k(m+M)g \cdot \varphi.$$ 
Over $N$ periods, the decrease in amplitude will be $4N$ times larger. Expressing the coefficient of rolling friction: $$k=\frac{m}{m+M} \frac{d \cdot \Delta\varphi_N}{4N}.$$

\begin{table}[H]
\centering
\begin{tabular}{|c|c|c|c|c|c|c|}
\hline
$x_{up}$, mm & $d$, mm & $\varphi_0$, $^\circ$ & $\varphi$, $^\circ$ & $N$ & $k$, mm & $\Delta k$, mm \\
\hline
0 & 186.5 & 30 & 10 & 18 & 0.23 & 0.019 \\
20 & 166.5 & 30 & 10 & 16 & 0.231 & 0.020 \\
40 & 146.5 & 30 & 10 & 14 & 0.233 & 0.022 \\
40 & 146.5 & 40 & 10 & 21 & 0.233 & 0.016 \\
40 & 146.5 & 40 & 20 & 13 & 0.25 & 0.025 \\
\hline
\end{tabular}
\end{table}

As a result, the rolling friction coefficient is $$k=(0.24 \pm 0.02)~\mathrm{mm}.$$

\end{solcard}

\begin{solcard}{A2}{Determine experimentally the sliding friction coefficient $\mu$ of a cardboard cylinder on a wooden block. Give the measurement method and installation diagram.}{4.00}

Using adhesive paste, attach the ruler as a rod to the base of the cylinder. Let's place this system on a wooden block. Let's start tilting the block (Fig. 1). At a certain angle of inclination $\beta$ of the block to the horizontal, the system will begin to slide. The friction coefficient is $$ \mu = \tan \beta = \tan(17^\circ \pm 1^\circ)= 0.31 \pm 0.02 .$$

\begin{figure}[H]
  \centering
  \includegraphics[width=0.72\textwidth]{E1_S_EN_assets/fig1.jpg}
\end{figure}
\end{solcard}

\end{document}